\section{One-sided short support and model evidence}\label{sec:models}
The complete Kloosterman matrix $\mathsf K=(S(m,n;q))_{1\le m,n<q}$ satisfies the directly verifiable identity
\[
\mathsf K\mathsf K^*=q^2I-(q+1)\mathbf J,
\]
where $\mathbf J$ is the all-ones matrix. To verify it, expand the two Kloosterman sums. Summation over $n$ gives $q\1_{r=s}-1$.
The first contribution is
$q\sum_{r\ne0}\e((m-m')r/q)=q(q\1_{m=m'}-1)$, and the second is $1$, giving the stated matrix. Every row sums to $1$.
Thus any vector $a$ with zero coordinate sum satisfies $\|\mathsf K^*a\|_2=q\|a\|_2$.
Choosing $a=e_{m_1}-e_{m_2}$ supported on two distinct indices shows that scale $q$ is still attained when one side has only two coordinates while the other is complete and centered.
One-sided short support therefore cannot replace the two-sided critical-array hypotheses.

This extremal example does not assert that the actual emitted coefficients attain that extremum.
It shows that a universal bound knowing only support and centering cannot be stronger.
Likewise, finite matrix checks test only the chosen scales. A model with a common growing-parameter theorem is stronger than a finite check, but still needs attachment of actual coefficients, a fixed starting point, parameter ranges, and residuals before it can address the original arithmetic problem.
We do not reissue historical finite certificates as asymptotic theorems or add synthetic experiments merely to increase the manuscript's length.
